\title{Parameter-Efficient Orthogonal Finetuning via Factorization}

\begin{abstract}
Large foundational models are increasingly prevalent, yet training them from the ground up remains prohibitively costly. Consequently, efficiently adapting these robust models to downstream tasks has become critically important. In this work, we investigate a well-founded finetuning strategy—Orthogonal Finetuning (OFT)—for downstream task adaptation. Although OFT exhibits strong generalization capabilities, it still requires a relatively large number of trainable parameters due to the high dimensionality inherent in orthogonal matrices. To overcome this, we first analyze OFT from the viewpoint of information transmission and identify several key criteria that promote improved parameter efficiency. Drawing inspiration from the Cooley-Tukey fast Fourier transform algorithm's ability to facilitate efficient information flow, we introduce a compact orthogonal parameterization based on butterfly structures. Incorporating this parameterization into OFT, we develop a novel parameter-efficient finetuning approach termed Orthogonal Butterfly (BOFT). BOFT generalizes the OFT framework, subsuming it as a special case, thereby offering a broader orthogonal finetuning methodology. Finally, we perform an extensive empirical evaluation on adapting large vision transformers, large language models, and text-to-image diffusion models across various vision and language downstream tasks.
\end{abstract}

\section{Introduction}

Recent models such as ChatGPT~\cite{brown2020language,chen2021evaluating} and Stable Diffusion~\cite{rombach2022high} showcase the impressive generalization capabilities of large foundational models. The rapid advancements in their performance have been accompanied by a substantial increase in parameter counts (for instance, GPT-3 contains roughly 175 billion parameters). This scale makes training foundational models from scratch increasingly impractical for researchers. Therefore, progress in the field hinges on the ability to adapt pretrained foundation models efficiently to downstream tasks without full retraining. Efficient task adaptation primarily falls into three categories: \emph{model finetuning}~\cite{hulora2022,zaken2022bitfit,guo2020parameter,raffel2020exploring,qiu2023controlling,zhang2022adaptive,chavan2023one}, which keeps the model architecture intact while finetuning a subset of parameters; \emph{adapter tuning}~\cite{rebuffi2017learning,houlsby2019parameter,pfeiffer2020adapterfusion,lin2020exploring,he2021towards}, which incorporates additional trainable modules into the original model; and \emph{prompt tuning}~\cite{li2021prefix,lester2021power}, which appends trainable prefix tokens to the input. Among these approaches, model finetuning stands out as a straightforward yet effective method that notably does not increase inference latency.

The core idea behind model finetuning is to maintain similarity between the pretrained and finetuned models according to specific metrics, thereby preserving the knowledge acquired during pretraining. For example, many existing finetuning techniques utilize a low learning rate as a heuristic to minimize the divergence between the pretrained and finetuned models. Considering the structured nature of weight matrices, a more principled strategy aims to retain the relational properties of these matrices, specifically the pairwise angles between neurons. This concept gives rise to a new finetuning framework called Orthogonal Finetuning~(OFT)~\cite{qiu2023controlling}, where neurons within the same layer are transformed using a common orthogonal matrix, ensuring that the pairwise angles between neurons are theoretically preserved throughout the finetuning process. While OFT has shown strong generalizability and stable convergence for finetuning text-to-image diffusion models~\cite{qiu2023controlling}, it can involve a large number of trainable parameters due to the high dimensionality of orthogonal matrices. To mitigate this, OFT employs a block-diagonal structure to reduce parameter count. Nonetheless, this gain in parameter efficiency comes with drawbacks—the orthogonal matrix follows a fixed sparsity pattern, and orthogonal transformations are applied independently across different blocks. This block-diagonal sparsity pattern, although parameter-saving, may impose unwanted inductive biases (e.g., limiting expressiveness and preventing approximation of classic linear transforms), and more critically, determining an effective sparsity pattern remains an open challenge.

The crucial challenge in tackling this issue is to produce a dense orthogonal matrix that remains parameter-efficient. Although it may initially appear impractical since a $d$-dimensional dense orthogonal matrix requires $\mathcal{O}(d^2)$ parameters, we propose an innovative strategy that constructs a dense orthogonal matrix by composing multiple factorized sparse matrices. This method is motivated by the observation that the number of trainable parameters can be decreased by exchanging computation time for storage space. Because the orthogonal matrix is represented as a product of sparse matrices, the multiplication must be executed repeatedly during each training iteration. To better understand this matrix factorization, we frame the generation of a dense orthogonal matrix as an information transmission problem. Within this framework, generating a dense orthogonal matrix via the multiplication of sparse matrices can be interpreted as transmitting information across a grid-structured graph. This perspective allows us to devise numerous sparse matrix factorization schemes that limit the number of trainable parameters while retaining sufficient expressiveness to generate dense matrices.

To attain parameter efficiency within our information transmission framework, we take inspiration from the butterfly structures found in the Cooley-Tukey fast Fourier transform algorithm~\cite{cooley1965algorithm}, where information is efficiently exchanged between nodes~\cite{klasing1994broadcasting}. In particular, the butterfly graph in the Cooley-Tukey algorithm naturally induces a sparse matrix factorization method known as butterfly factorization. Using butterfly factorization, we can produce a $d$-dimensional dense matrix as a product of $\mathcal{O}(\log d)$ sparse matrices, each containing $\mathcal{O}(d)$ non-zero elements. By ensuring the orthogonality of each sparse matrix, we obtain an efficient orthogonal parameterization with only $\mathcal{O}(d \log d)$ parameters, representing a substantial reduction compared to the original parameter count. Exploiting this efficient orthogonal parameterization, we introduce a novel parameter-efficient finetuning technique called Orthogonal Butterfly~(BOFT). BOFT generalizes OFT as a special case and establishes a broad orthogonal finetuning framework. A shared attribute between the block-diagonal and butterfly structures is sparsity: both impose data sparsity on orthogonal matrices to decrease the effective number of trainable parameters. It is noteworthy to contrast our approach with the low-rank structure used in LoRA~\cite{hulora2022}; low-rank matrices and sparse matrices are two primary classes of structured matrices~\cite{candes2011robust} that facilitate parameter efficiency.

In contrast to the block-diagonal structure employed by OFT to balance expressiveness and regularity, BOFT utilizes the butterfly structure to create a more gradual interpolation between matrices belonging to the full orthogonal group (expressiveness) and identity matrices (regularity). This approach allows us to identify a smaller hypothesis space within the orthogonal group tailored for downstream tasks. Given that the butterfly structure is extensively used in numerous fast linear transformations, such as the discrete Fourier and discrete cosine transforms, we contend that our structured strategy for parameter efficiency introduces an implicit inductive bias that enhances generalization and mitigates overfitting. Our reasoning is based on the fact that the butterfly structure can readily reconstruct many classical linear transforms, whereas the block-diagonal structure in OFT cannot recover any. Our key contributions are summarized as follows:

\begin{itemize}[leftmargin=*,nosep]
    \item We address the challenge of parameter efficiency in orthogonal finetuning through a novel information transmission framework. This framework reformulates the task of designing a parameter-efficient dense orthogonal matrix as an information transmission problem over a grid-structured graph.
    \item Drawing inspiration from the butterfly structures found in the Cooley-Tukey algorithm, we introduce Orthogonal Butterfly, a parameter-efficient orthogonal finetuning technique, developed within the information transmission framework.
    \item We offer several theoretical insights explaining why BOFT can substantially decrease the number of trainable parameters while maintaining strong expressiveness and training stability. Moreover, due to matrix factorization, BOFT features an interesting weight interpolation property (refer to Figure~\ref{fig:boft_interpolation}).
    \item For the first time, we extend orthogonal finetuning to a range of tasks beyond controllable text-to-image generation~\cite{qiu2023controlling}, demonstrating its broad potential as a universal model finetuning approach. Specifically, we apply BOFT to diverse adaptation scenarios spanning computer vision and natural language processing. BOFT consistently surpasses current state-of-the-art methods by a significant margin, confirming its superior parameter efficiency and generalization capabilities.
\end{itemize}

\section{Related Work}

\textbf{Parameter-efficient finetuning (PEFT)}. As foundational models (\emph{e.g.}, \cite{brown2020language,radford2021learning,rombach2022high,kirillov2023segment}) continue to grow larger and more capable, the interest in finetuning these models for downstream tasks in a parameter-efficient manner has increased significantly~\cite{treviso2023efficient,ding2023parameter,lialin2023scaling}. Among the numerous PEFT techniques~\cite{houlsby2019parameter,aghajanyan2020intrinsic,hulora2022,edalati2022krona,wang2022adamix,gheini2021cross,zaken2022bitfit,guo2020parameter,sung2021training,ansell2022composable,lester2021power,li2021prefix,vu2022spot,he2021towards,mao2021unipelt,karimi2021compacter,liu2022few,sung2022lst,chen2022parameter,jia2022visual,chen2022adaptformer,zhang2022neural,jie2023fact,lian2022scaling,luo2023towards}, reparameterization-based approaches~\cite{aghajanyan2020intrinsic,hulora2022,edalati2022krona} are particularly relevant to our study. LoRA~\cite{hulora2022} modifies the pretrained weight matrix by adding the product of two low-rank matrices, delivering strong performance on natural language tasks. Given that LoRA fixes the rank of these added matrices, some methods \cite{zhang2022adaptive,valipour2022dylora,zhang2023increlora} adaptively vary the rank across different layers to better distribute the parameter budget. Owing to its straightforwardness, this low-rank weight reparameterization technique has become widely adopted \cite{dettmers2023qlora,zi2023delta,chavan2023one}. Drawing inspiration from the role of hyperspherical energy in characterizing generalization~\cite{liu2018learning,liu2021orthogonal}, \cite{qiu2023controlling} introduces orthogonal finetuning, an alternative yet effective approach for finetuning text-to-image diffusion models. Specifically, OFT learns an orthogonal matrix that transforms neurons within the same layer, achieving enhanced generalization and consistently more stable training compared to LoRA. Despite its strong results, OFT typically involves a larger number of trainable parameters than LoRA. Consequently, improving the parameter efficiency of OFT is a worthwhile objective. Furthermore, it remains unclear whether OFT can be extended to a broader range of adaptation tasks beyond controlling text-to-image diffusion models~\cite{qiu2023controlling}. BOFT enhances OFT’s parameter efficiency by leveraging butterfly factorization. This advancement enables us to showcase the effectiveness of orthogonal finetuning across general adaptation tasks.

\textbf{Butterfly structures}. The radix-2 Cooley-Tukey algorithm recursively decomposes an $N$-point discrete Fourier transform into two $\frac{N}{2}$-point discrete Fourier transforms, producing a butterfly structure representable as a product of multiple sparse matrices (commonly referred to as a butterfly matrix). Butterfly matrices have been employed to parameterize orthogonal matrices, avoiding pivoting in Gaussian elimination and boosting efficiency~\cite{parker1995random}, to stabilize recurrent neural network training~\cite{jing2017tunable}, and in kernel approximation~\cite{munkhoeva2018quadrature}. Additionally, \cite{dao2019learning,dao2019kaleidoscope} learn fast linear transforms using butterfly parameterizations. Works such as \cite{chen2022pixelated,dao2022monarch} apply butterfly matrices to facilitate efficient neural network training. Butterfly structures are also valuable in fast matrix-vector multiplication~\cite{michielssen1996multilevel,o2010algorithm}, data-sparse matrix approximation~\cite{li2015butterfly}, and network transmission~\cite{klasing1994broadcasting,li2003linear,rajasingh2016transmission}. In contrast to prior studies, our focus is on exploiting butterfly structures to improve the parameter efficiency of OFT.

\section{An Information Transmission View on Orthogonal Finetuning}

We begin with some foundational aspects of OFT. To finetune the pretrained weight matrix, OFT reparameterizes the updated weight matrix as the product of a learnable orthogonal matrix and the fixed pretrained weight matrix. Unlike LoRA, which updates weights by adding a low-rank matrix, OFT employs a multiplicative orthogonal matrix for weight updates. To enhance parameter efficiency, LoRA leverages a low-rank structure, whereas the original OFT adopts a block-diagonal orthogonal structure~\cite{qiu2023controlling} (smaller diagonal block sizes lead to greater parameter efficiency in OFT). An intuitive comparison is shown in Figure~\ref{fig:comp_lora}. The rationale for using orthogonal transformations in finetuning the weight matrix is to maintain the pairwise angles between neurons~\cite{liu2018learning,liu2021orthogonal,qiu2023controlling}, thereby preserving much of the semantic knowledge acquired during pretraining. Specifically, OFT optimizes an orthogonal matrix $\bm{R}\in\mathbb{R}^{d\times d}$ for a pretrained linear layer $\bm{W}^0\in\mathbb{R}^{d\times n}$, modifying the forward computation from $\bm{z}=(\bm{W}^0)^\top\bm{x}$ to $\bm{z}=(\bm{R}\bm{W}^0)^\top\bm{x}$, where $\bm{x}\in\mathbb{R}^{d}$ and $\bm{z}\in\mathbb{R}^{n}$ denote the input and output vectors, respectively. To ensure zero initialization, $\bm{R}$ is initialized as the identity matrix. To maintain the orthogonality of $\bm{R}$ throughout finetuning, we adopt the Cayley parameterization as in \cite{qiu2023controlling,liu2021orthogonal}, i.e., $\bm{R}=(\bm{I}+\bm{Q})(\bm{I}-\bm{Q})^{-1}$, where $\bm{Q}$ is a skew-symmetric matrix satisfying $\bm{Q}=-\bm{Q}^\top$. For parameter efficiency, the block-diagonal structure is enforced by representing $\bm{R}$ as $\text{diag}(\bm{R}_1,\bm{R}_2,\cdots,\bm{R}_r)$, with each $\bm{R}_i\in\mathbb{R}^{b\times b}$ being a small orthogonal matrix and $br=d$. The parameter efficiency gained from the block-diagonal structure comes with a trade-off — it assumes that the dimensions of a neuron (i.e., a column vector of the weight matrix $\bm{W}^0$) are partitioned into $r$ groups, each transformed independently by different orthogonal matrices. Although this block-diagonal structure proves empirically effective, splitting the dimensions of a neuron based solely on their indices lacks intuitive justification, motivating the preference for dense orthogonal matrices. This raises a natural question: \emph{Is it possible to construct a dense orthogonal matrix without compromising parameter efficiency?}

To tackle this issue, we suggest constructing a dense orthogonal matrix as a product of several sparse orthogonal matrices. To offer a coherent and intuitive framework for examining the sparsity pattern in orthogonal matrix factorization, we reinterpret the task of generating a dense orthogonal matrix in OFT as an information transmission problem. More precisely, creating a dense matrix $\bm{R}\in\mathbb{R}^{d\times d}$ by multiplying $m$ square matrices, i.e., $\bm{R}=\bm{B}_m\bm{B}_{m-1}\cdots\bm{B}_1$, can be viewed as transmitting information across a grid composed of $d\times (m+1)$ nodes, as depicted in Figure~\ref{fig:info_trans}. The rationale for this information transmission perspective stems from the insight that a $d$-dimensional dense square matrix can be understood as representing dense connectivity from $d$ nodes to another set of $d$ nodes. For the matrix $\bm{R}$, a zero element $R_{ij}$ implies that information cannot flow from the $j$-th node to the $i$-th node, while a non-zero element indicates that information transfer is possible. Consequently, expressing the dense matrix $\bm{R}$ as a product of matrices $\bm{B}_m\bm{B}_{m-1}\cdots\bm{B}_1$ can be interpreted as performing a sequential information exchange governed by the graphs induced by each $\bm{B}_i, \forall i$. The information flows initially through $\bm{B}_1$ and finally through $\bm{B}_m$. 

To illustrate, consider the factorization $\bm{R}=\bm{B}_5\bm{B}_4\bm{B}_3\bm{B}_2\bm{B}_1$ shown in Figure~\ref{fig:info_trans}, where both sparsity patterns and the induced graph are presented. The graph is derived from unrolling the matrix multiplication. In this induced graph, each matrix $\bm{B}_i$ represents the connectivity between the nodes at level $i$ and those at level $i+1$. More precisely, the element at position $(j_1,j_2)$ in $\bm{B}_i$ indicates the presence or absence of a directed edge from the $j_2$-th node at level $i$ to the $j_1$-th node at level $i+1$ (with zero meaning no edge). For the product $\bm{B}_5\bm{B}_4\bm{B}_3\bm{B}_2\bm{B}_1$ to form a dense matrix, it is necessary that every node at the first level can send information to all nodes at the sixth level. When considering only $\bm{R}=\bm{B}_4\bm{B}_3\bm{B}_2\bm{B}_1$, linking source nodes at the first level to receiver nodes at the fifth level, it is evident that information from node $1$ cannot reach node $3$. Thus, the sparsity pattern of $\bm{B}_4\bm{B}_3\bm{B}_2\bm{B}_1$ features a zero at the $(3,1)$ position. This same relationship between the induced graph and sparsity pattern holds for the factorizations $\bm{R}=\bm{B}_3\bm{B}_2\bm{B}_1$ and $\bm{R}=\bm{B}_2\bm{B}_1$. More broadly, for a matrix $\bm{R}\in\mathbb{R}^{d\times d}$ to be dense, the factorization matrices $\bm{B}_m,\ldots,\bm{B}_1$ must correspond to a set of directed edges on a $d\times (m+1)$ grid, where each directed edge connects only nodes in adjacent levels (i.e., adjacent columns), ensuring that information from every node at the first level is capable of reaching every node at the $(m+1)$-th level.

The perspective of information transmission aids in deepening our comprehension of the sparsity pattern observed in the factorization matrices within OFT. Figure~\ref{fig:blk_diag} illustrates the block-diagonal form of $\bm{R}$ in the original OFT. Although this block-diagonal configuration decreases the number of trainable parameters, it cannot produce a dense matrix $\bm{R}$. Our objective is to construct a dense orthogonal matrix by combining $m$ sparse orthogonal matrices while utilizing as few effective trainable parameters as possible. From the information transmission perspective, the primary criteria for achieving this goal are \emph{(i)} \textbf{dense connectivity}: ensuring each node in the initial layer connects to at least one node in the final layer, and \emph{(ii)} \textbf{minimum free edges}: minimizing the total number of edges subject to the orthogonality constraint. It is important to highlight that orthogonality imposes a subtle constraint on the edges connecting adjacent layers. For instance, for each matrix $\bm{B}_i$ to be full-rank (a necessary condition for orthogonality), there must be at least $d$ edges creating a one-to-one correspondence between all nodes in the $i$-th layer and those in the $(i+1)$-th layer. This implies that the number of edges between adjacent layers is at least $d$ (\emph{e.g.}, 4 in the example shown in Figure~\ref{fig:info_trans}). These $d$ edges are essential for orthogonality and should be excluded from the count of edges because these elements are non-trainable (\emph{e.g.}, in a $d\times d$ orthogonal matrix with $d$ non-zero entries, these entries can only take values of $\pm 1$). Since orthogonal matrices require fewer parameters than full matrices, the orthogonality constraint introduces additional dependencies among edges. For example, in a $2\times 2$ orthogonal matrix, having a zero at position $(1,1)$ necessarily implies a zero at position $(2,2)$ (\emph{i.e.}, the absence of one edge implies the absence of another). Therefore, for each valid set of edge connections, orthogonality may impose the addition or removal of some edges. By examining the non-zero structure of the composed orthogonal matrix, the information transmission framework proves especially valuable in OFT, as our focus is solely on the non-zero trainable elements of $\bm{R}$, regardless of their specific values.

A straightforward dense connection between two layers requires $\mathcal{O}(d^2)$ edges (\emph{i.e.}, one dense orthogonal matrix), resulting in $d^2-d$ edges (for $d=4$, this equals 12 edges). Figure~\ref{fig:info_trans} presents an instance of a valid matrix factorization that uses a total of $10$ edges, which is actually fewer than a single dense orthogonal matrix. This approach allows us to analyze the parameter efficiency of OFT through a graphical lens, making it simple to derive feasible factorizations within this framework. We take inspiration from a fascinating topology in the Cooley-Tukey algorithm, known as butterfly graphs~\cite{cooley1965algorithm}, which can connect $d$ source nodes to $d$ receiver nodes densely and efficiently with $\mathcal{O}(d\log d)$ edges. For instance, the structure shown in Figure~\ref{fig:info_trans} uses $10$ edges to realize dense connectivity, whereas the butterfly network requires only $8$ edges. Next, we explain how the butterfly architecture can enhance parameter efficiency.

\section{Orthogonal Parameterization by Butterfly Factorization}

The butterfly architecture was initially introduced in the Cooley-Tukey algorithm to facilitate fast Fourier transform. In the Fourier transform, a local modification in the frequency domain can trigger a global variation in the spatial domain, which is conceptually analogous to our information transmission problem—each node in the first layer can send information to every node in the last layer. The butterfly structure has also become a widely used computer network topology~\cite{solihin2015fundamentals,li2011linear}, known for enabling efficient information exchange. Assuming $k \geq 2$ is a power of $2$, we first define the \emph{butterfly factor} $\bm{B}^F(k)$ as

\begin{equation}\label{eq:but_factor}
\begin{aligned}
    \bm{B}^F(k)=\begin{bmatrix}
    \text{diag}(\bm{d}_1) & \text{diag}(\bm{d}_2) \\
    \text{diag}(\bm{d}_3) & \text{diag}(\bm{d}_4)
    \end{bmatrix}\in\mathbb{R}^{k\times k},
\end{aligned}
\end{equation}

where each $\bm{d}_i\in\mathbb{R}^{\frac{k}{2}}, \forall i$ denotes a vector. For $d=2^N$, we recursively define the $d$-dimensional \emph{butterfly matrix} $\bm{B}(d)\in\mathbb{R}^{d\times d}$ as

\begin{equation}
\begin{aligned}
    \!\!\bm{B}(d)=\tilde{\bm{B}}(d,d)\!\cdot\!\begin{bmatrix}
    \bm{B}_1(\frac{d}{2})\!\!\!\!\! & \bm{0} \\
    \bm{0}\!\!\!\!\! &\bm{B}_2(\frac{d}{2})
    \end{bmatrix}= \tilde{\bm{B}}(d,d)\tilde{\bm{B}}(d,\frac{d}{2})\cdots\tilde{\bm{B}}(d,2),
\end{aligned}
\end{equation}

where $\bm{B}_1(\frac{d}{2})$ and $\bm{B}_2(\frac{d}{2})$ denote two butterfly matrices of dimension $\frac{d}{2}$. Next, we define the \emph{butterfly component} as $\thickmuskip=2mu \medmuskip=2mu \tilde{\bm{B}}(d,k)=\text{diag}(\bm{B}^F_1(k),\cdots,\bm{B}^F_{d/k}(k))$, which is a block-diagonal matrix of size $d \times d$ composed of blocks of size $k$. Here, $\bm{B}^F_i(k), \forall i$ represent butterfly factors introduced in Equation~\ref{eq:but_factor}. We are now prepared to use the butterfly matrix for parameterizing an orthogonal matrix. To do this, it suffices to ensure that every multiplicative factor $\tilde{\bm{B}}(d,k), \forall k$ within the butterfly matrix $\bm{B}(d)$ is orthogonal. We begin by examining the block-diagonal matrix $\tilde{\bm{B}}(d,2)$ with block size 2; it is straightforward to guarantee orthogonality of $\tilde{\bm{B}}(d,2)$ by employing the Cayley transform (or equivalently a 2-dimensional rotation) to parameterize each block, analogous to \cite{liu2021orthogonal,qiu2023controlling}. The pattern of non-zero entries in each butterfly component can be interpreted as a permutation of the non-zero structure of $\tilde{\bm{B}}(d,2)$, which implies all butterfly components can be efficiently parameterized as orthogonal matrices. This framework thus provides an effective parameterization of orthogonal matrices constructed from numerous $2 \times 2$ orthogonal blocks. Extending the butterfly matrices as in \cite{chen2022pixelated}, we define a block butterfly component $\tilde{\bm{B}}^b(d,k)$ in which every entry of $\bm{d}_i, \forall i$ is replaced by a $b \times b$ matrix. To ensure that the block butterfly component $\tilde{\bm{B}}^b(d,2)$ is orthogonal, we parameterize each $2b \times 2b$ block to be orthogonal. The sparsity pattern of other butterfly components $\tilde{\bm{B}}^b(d,k)$ with $k > 2$ corresponds to block-wise permutations of the non-zero pattern of $\tilde{\bm{B}}^b(d,2)$ and thus can be similarly parameterized to form orthogonal matrices. Putting all elements together, the forward pass in BOFT is

\begin{equation}
    \begin{aligned}\nonumber
    \bm{z}=\big(\bm{R}(m,b)\cdot\bm{W}^0\big)^\top\bm{x},~~~~\text{s.t.}~\bigg{\{}\bm{R}(m,b)=\prod_{i=1}^m \tilde{\bm{B}}^b_{(d,i)}~~\&~~\underbrace{\big(\tilde{\bm{B}}^b_{(d,j)}\big)^\top\tilde{\bm{B}}^b_{(d,j)}=\tilde{\bm{B}}^b_{(d,j)}\big(\tilde{\bm{B}}^b_{(d,j)}\big)^\top\!=\bm{I}_d}_{\forall j\in [1, m]}\bigg{\}},
    \end{aligned}
\end{equation}

Here, for simplicity, we denote $\tilde{\bm{B}}^b(d,2^{m - i+1})$ as $\tilde{\bm{B}}^b_{(d,i)}$, and $\bm{I}_d$ represents the identity matrix of dimension $d$. The orthogonal matrix $\bm{R}(m,b)\in\mathbb{R}^{d\times d}$ is constructed as the product of several orthogonal butterfly components. For ease of notation, we define BOFT with $\bm{R}(m,\frac{b}{2})$ as BOFT($m$,$b$), where $b \geq 2$. When $m=1$, BOFT($1$,$b$) corresponds to the block-diagonal OFT~\cite{qiu2023controlling} with block size $b$. BOFT($1$,$d$) is equivalent to the original OFT~\cite{qiu2023controlling} with a fully unconstrained orthogonal matrix. BOFT($\log \frac{2d}{b}$,$b$) employs the block butterfly matrix $\bm{B}^{\frac{b}{2}}(d)$ as $\bm{R}$, producing a dense orthogonal matrix $\bm{R}$. Generally, BOFT($m$,$b$) involves $\frac{1}{2}(b-1)dm$ effective trainable parameters when fine-tuning a linear layer sized $d \times n$. Using the butterfly matrix, i.e., $m=\log d$ and $b=2$, BOFT requires $\mathcal{O}(d \log d)$ parameters. By comparison, the original OFT with a fully dense orthogonal matrix uses $\mathcal{O}(d^2)$ parameters, while the block-diagonal OFT with block count $r$ utilizes $\mathcal{O}(bd)$. Hence, the original OFT must have block size $b=d$ to create a dense orthogonal matrix, whereas BOFT can achieve this with any choice of $b$.

\textbf{Identity initialization for BOFT}. Fine-tuning approaches typically commence from the exact pretrained model so that the fine-tuned version does not diverge significantly from the original. For instance, LoRA initializes low-rank weights to zero. In BOFT, we initialize all butterfly components as identity matrices (i.e., the skew-symmetric matrices used in the Cayley transform are initialized to zero).

\textbf{Multiplicative Dropout for BOFT}. LoRA~\cite{hulora2022} additionally applies a Dropout layer on the low-rank weight updates to mitigate overfitting. While the standard Dropout~\cite{srivastava2014dropout} naturally suits LoRA, it is not compatible with BOFT due to the multiplicative nature of the weight update. To overcome this, we introduce a multiplicative Dropout tailored for BOFT. Since the orthogonal matrix $\bm{R}(m,b)$ consists of $m$ orthogonal butterfly components, which can be rearranged into $2b \times 2b$ block-diagonal orthogonal matrices, the multiplicative Dropout randomly selects $p_1$ percent of these butterfly components and $p_2$ percent of the diagonal blocks within each component, subsequently replacing those blocks with identity matrices.

\section{Intriguing Insights and Discussions}

\textbf{Expressivity of BOFT}. The butterfly structure combined with permutations can exactly recover many well-known fast linear transforms~\cite{dao2019learning,dao2019kaleidoscope} (such as the fast Fourier transform and Hadamard transform); however, the extent to which our orthogonal butterfly matrix can approximate a general orthogonal matrix is still unclear. We begin by performing a simulation to approximate a random dense orthogonal matrix~\cite{anderson1987generation} of size $512\times 512$. The outcomes shown in Figure~\ref{fig:para_eff} are averaged over 10 different random seeds. The y-axis represents the approximation error, while the x-axis indicates the number of effective trainable parameters. Each curve of the same color corresponds to BOFT with an identical block size, and the leftmost point indicates the error for BOFT($1$,$b$) (i.e., the original block-diagonal OFT with block size $b$). Overall, BOFT tends to offer superior parameter efficiency compared to OFT. For instance, BOFT($9$,$2$) demonstrates better expressiveness than BOFT($1$,$16$) but requires substantially fewer parameters. Typically, BOFT with a smaller $b$ and larger $m$ is more parameter-efficient. For example, BOFT($6$,$4$) uses considerably fewer parameters while achieving an approximation error comparable to BOFT($2$,$16$). In general, the butterfly matrix embodies a more structured subset of the orthogonal group (relative to the block-diagonal structure), making BOFT provably more expressive than OFT given the same block size.

\begin{theorem}[Expressivity of BOFT]\label{thm:exp_boft}
    BOFT exhibits greater expressiveness than OFT when using the same block size. To approximate every orthogonal matrix of size $d$ with a butterfly matrix, one can multiply butterfly matrices in the form $\bm{B}_{d-1,1}(d)\bm{B}^\top_{d-1,2}(d)\cdots\bm{B}_{1,1}(d)\bm{B}^\top_{1,2}(d)$, where $\bm{B}_{i,j}(d)$ are butterfly matrices for all $i$ and $j$.
\end{theorem}

Theorem~\ref{thm:exp_boft} indicates a straightforward extension for BOFT — the final orthogonal matrix can be generalized as $\thickmuskip=2mu \medmuskip=2mu \bm{R}^G(m_1,b_1,m_2,b_2,l)=\bm{R}_{l,1}(m_1,b_1)\bm{R}_{l,2}^T(m_2,b_2)\cdots\bm{R}_{1,1}(m_1,b_1)\bm{R}_{1,2}^T(m_2,b_2)$, where $\bm{R}_{i,j}^T(m,b)$ represents the orthogonal matrix used in BOFT. When $m_1=m_2=\log d$, $b_1=b_2=2$, and $l=d-1$, the matrix $\bm{R}^G(m_1,b_1,m_2,b_2,l)$ can represent the entire orthogonal group. This type of matrix composition is also referred to as the kaleidoscope hierarchy~\cite{dao2019kaleidoscope}. Nonetheless, we emphasize that increased expressiveness does not necessarily translate into improved performance during finetuning, since full finetuning, although universally expressive, frequently results in unsatisfactory outcomes. Balancing expressivity and regularity is crucial for ensuring the generalizability of model finetuning. BOFT expands the finetuning parameter space by incorporating structural priors, which allows us to achieve a better balance between expressivity and regularity.

\textbf{Spectral properties}. Orthogonal finetuning typically results in superior spectral characteristics compared to LoRA, as it exactly maintains the spectral norm of the pretrained weight matrix $\bm{W}^0$. This can be understood through the singular value decomposition: $\bm{W}^0=\bm{U}\bm{\Sigma}\bm{V}^\top$, where $\bm{U}$ and $\bm{V}$ are orthogonal matrices, and $\bm{\Sigma}$ is a diagonal matrix of singular values. Both OFT and BOFT apply an orthogonal matrix $\bm{R}$ from the left, producing the finetuned weights $\bm{R}\bm{U}\bm{\Sigma}\bm{V}^\top$, which does not alter the largest singular value (i.e., the spectral norm of $\bm{W}^0$). Preserving this property has been demonstrated to significantly enhance training stability and generalization~\cite{miyato2018spectral,yoshida2017spectral}. Additional intriguing mathematical properties are discussed in Appendix~\ref{app:sn_property}.

\textbf{Orthogonal finetuning as learning bilinear similarity}. BOFT can be expressed as learning the bilinear similarity $\bm{w}^0_i\bm{R}\bm{x}$, where $\bm{w}^0_i$ denotes the $i$-th neuron (i.e., column vector) of the weight matrix $\bm{W}^0$. In this view, BOFT corresponds to learning the bilinear similarity matrix $\bm{R}$ subject to a strong constraint—namely, that $\bm{R}$ is orthogonal—thereby inherently linking it to distance metric learning~\cite{xing2002distance} and bilinear forms~\cite{roman2005advanced}.

\textbf{Inductive bias and generalization in BOFT}. Because $\bm{R}(m,b)$ in BOFT generally belongs to a structured subset of the orthogonal group that restricts the hypothesis space, BOFT naturally introduces an inductive bias. We contend that the structured inductive bias stemming from butterfly factorization favors generalization, as it embodies shared structural patterns present in numerous classical linear transforms~\cite{dao2019kaleidoscope}, including discrete Fourier transform, discrete sine/cosine transform, and Hadamard transform. Furthermore, the sparse matrix factorization inherent in BOFT may confer additional implicit inductive biases~\cite{gunasekar2017implicit,li2018algorithmic,liu2019neural}.

\textbf{Comparison to butterfly-based sparse training}. Several studies~\cite{dao2019learning,dao2019kaleidoscope,chen2022pixelated,dao2022monarch} explore sparse training utilizing the butterfly parameterization. These works generally concentrate on directly reparameterizing weight matrices using the butterfly structure and training neural networks from the beginning. \cite{dao2022monarch} investigates finetuning pretrained weights by initially projecting them onto a variant of butterfly matrices, followed by optimizing the projected components for downstream applications. In contrast, BOFT introduces a distinct finetuning approach that modifies the weights through layer-shared weight matrices.

\input{rebuttal-results}

\section{Concluding Remarks and Limitations}

In this paper, we present Orthogonal Butterfly (BOFT), a universal parameter-efficient finetuning technique for foundation models built upon the butterfly architecture. The central idea to enhance parameter efficiency is to represent a dense orthogonal matrix as a product of several sparse orthogonal matrices. To facilitate the discovery of suitable matrix factorizations, we introduce a graphical information transmission framework. Utilizing this framework, we demonstrate that the butterfly structure effectively satisfies our criteria for sparse orthogonal matrix factorization. We empirically confirm the effectiveness of BOFT in finetuning large language models, extensive vision models, and text-to-image generative models. Our experiments further corroborate the advantages of BOFT as a versatile finetuning method.

Although BOFT shows strong empirical performance, it is not without limitations. Because the ultimate orthogonal matrix in BOFT results from multiplying multiple orthogonal matrices, the training runtime overhead is somewhat higher compared to OFT. Addressing how to reduce BOFT's training runtime is an open challenge. Fortunately, once finetuning concludes, the orthogonal matrices learned via BOFT can be integrated directly into the pretrained model without causing additional inference latency. Additionally, it remains unclear if the butterfly network represents the optimal structure for information transmission. Our information transmission framework allows us to draw from insights in a different field — computer networking — where the efficiency of network topologies for information flow is extensively studied. We anticipate that more efficient network architectures could be leveraged to construct dense orthogonal matrices.

\end{document}